\GuideSpread{The shared prime is the conjecture's conclusion}{An exact positive example with all exponents greater than two}{01_beal_terms.pdf}{
Beal asks whether every equality $A^x+B^y=C^z$ in positive integers, with $x,y,z>2$, forces a prime dividing all three bases.\par
Here $3^3+6^3=3^5$ is the equality $27+216=243$. Its bases have common divisor three, so it illustrates the proposed conclusion. It is not a counterexample.\par
If a prime divides two bases in such an equality, it divides the third. Thus a hypothetical counterexample has pairwise coprime bases. This primitive condition belongs to the bases, rather than to an arbitrary rational parametrization. The familiar $3^2+4^2=5^2$ lies outside the conjecture because its exponents are two.
}{Keep the exponent threshold and the common-prime condition visible together.}

\GuideSpread{One signature is only one branch of Beal}{Exponent absorption retains odd primes and exponent four}{02_exponent_reduction.pdf}{
When an odd prime $p$ divides $x$, write $A^x=(A^{x/p})^p$. If $x$ has no odd prime divisor, then $x>2$ is divisible by four, giving $A^x=(A^{x/4})^4$. The same operation applies to the other bases.\par
Raising a base to a positive integer power preserves its prime support. Primitive bases therefore remain primitive after exponent absorption.\par
The figure applies this elementary reduction to the illustrative signature $(12,15,14)$, producing $(4,3,7)$. It does not produce an equation or solution. The reduced family remains infinite, since odd primes remain unrestricted. The focused signature $(3,5,7)$ is one substantial branch, rather than the full conjecture.
}{A complete proof for one signature needs a separate argument before it can cover every signature.}

\GuideSpread{Three valuation sectors recognize primitive bases}{The focused signature $A^3+B^5=C^7$}{03_valuation_packets.pdf}{
Set $t=A^3/C^7$ and $s=1-t=B^5/C^7$. For primitive positive integer bases, $0<t<1$. At an active numerator prime of $t$, its valuation is a positive multiple of three. At an active numerator prime of $s$, its valuation is a positive multiple of five.\par
At a denominator prime, both valuations are equal negative multiples of seven. A fourth, inactive sector has both valuations zero.\par
Conversely, the complete valuation criterion at every rational prime gives a unique positive rational scaling that recovers primitive integer bases from the universal parametrization on the next page. This is an exact recognition equivalence. It does not guarantee that a suitable rational $t$ exists.
}{The three distinct exponent divisibilities must hold for the same rational parameter.}

\GuideSpread{A rational identity keeps its integrality problem}{The positive universal parametrization}{04_rational_scaling.pdf}{
For $s=1-t$, the displayed powers satisfy $A_0^3+B_0^5=C_0^7$ identically. All three terms have common factor $t^{35}s^{84}$, and the remaining sum is $t+s=1$.\par
Weighted scaling by one positive rational $\lambda$ preserves the equation, since $3\cdot35=5\cdot21=7\cdot15=105$.\par
At a prime in the first, second, or denominator sector, the forced valuation of $\lambda$ is respectively $-v(t)/3$, $-4v(s)/5$, or $8[-v(t)]/7$. Those valuations must be integers. This recovers the preceding divisibility test. A rational formula that makes the equation automatic has not made primitive integrality automatic.
}{Preserve the scaling valuation and the three simultaneous power conditions.}

\GuideSpread{The local coefficient sieve has explicit survivors}{The squarefree stratum $(a,b,c)=(3,5,2)$}{05_local_saturation.pdf}{
Writing $A=au^2$, $B=bv^2$, and $C=cw^2$ gives $a^3u^6+b^5v^{10}=c^7w^{14}$. For the displayed coefficients, the dyadic test is $a+b\equiv0\pmod8$. At the odd coefficient primes, the required quadratic residues are $bc=10$ at three and $ac=6$ at five.\par
These tests pass. The written local criterion and lifting argument provide a primitive solution over every $\Z_p$, with separate local roots. Positive real solvability also holds.\par
This is not one rational or integer solution. Refining the same projected coefficient conditions to higher prime powers cannot eliminate this stratum. A global method must couple the local data through additional arithmetic structure.
}{Placewise solvability and a single global equality are different mathematical assertions.}

\GuideSpread{The base field and coefficient field must stay distinct}{The two appearances of the symbol five}{06_two_fields.pdf}{
The number field is $K=\Q(\sqrt5)$, with fundamental unit $\phi=(1+\sqrt5)/2$. It describes arithmetic places, global units, and ramification.\par
The residual coefficient field is $k=\F_{49}$. Its presentation is $\F_7[w]/(w^2-5)$. Since five is a nonsquare modulo seven, this is a degree-two field over $\F_7$. The code $a+7b$ stores the coefficient $a+bw$ with $0\le a,b<7$.\par
The radicals have different roles and different characteristics. An arithmetic descent involving the automorphism of $K/\Q$ requires its own compatibility argument. It is not licensed by using similar notation for coefficient conjugation. All later finite matrices are over $\F_7$, intentionally extended to $k$.
}{Always state which field contains a coefficient, a unit, a prime, or a radical.}

\GuideSpread{Joint trace labels retain exact information}{A source-linked example at the auxiliary prime eleven}{07_trace_correlation.pdf}{
The project source data give fifteen regular rank-two trace residues in $k$ at eleven. The guide recomputes these roots from the stored exact trace polynomials. The paired traces are $(a,a^7)$, retaining coefficient conjugation.\par
The dark points are the fifteen admissible conjugate pairs. Taking the two marginal supports independently gives all $15^2=225$ pale positions. That operation adds 210 pairs that violate the retained correlation.\par
Axis positions are ordered trace labels, not real values of finite-field elements. The example demonstrates information loss from independent projection. It does not establish that any discarded pair occurs for a surviving modular form, or that this correlation alone excludes a branch.
}{A selector should preserve the joint label when the arithmetic object preserves it.}

\GuideSpread{Precision belongs to the same natural frame}{Residual data, first digits, and further integral constraints}{08_depth_frame.pdf}{
The current bounded residual patterns have lattice depth one. The natural endpoint supplies the residual class $\alpha$. The opposite endpoint supplies $\beta$, which appears as a scaled first digit in the common natural lattice modulo $49$. The first-stage cochain is $p$, while the upper digit is $q$.\par
An integral representation includes higher compatibility conditions, including the next precision modulo $343$. A matrix satisfying a first-order equation need not lift to the actual integral Frey object.\par
The route uses one pinned natural lattice, common frame, and reference $\rho_0$. Changing frames independently at each place can alter the apparent local target without preserving its global meaning. General depth arguments in the research record retain their own precision and hypotheses. They do not reuse the depth-one table as a higher-depth table.
}{Track the frame and reference with the precision of every claim.}

\GuideSpread{A local packet obeys a group relation}{Compatible tame inertia $T$ and arithmetic Frobenius $F$}{09_tame_packet.pdf}{
At a depth-one beta-supporting place, the first-order toric packet has $\rho(T)=I+7M$ modulo $49$, with $M=c\bigl(\begin{smallmatrix}\theta&-\theta^2\\1&-\theta\end{smallmatrix}\bigr)$ and $c=\beta(T)\ne0$. Its square is zero over $k$.\par
Write $b=\beta(F)$ and $h=(Q-1)/7\bmod7$, where $Q$ is the residue norm. The packet gives $p(T)=\theta c$ and $p(F)=\theta b-h$. Eliminating $\theta$ yields $cp(F)-bp(T)=-hc$.\par
The tame relation and this affine equation concern the same compatible local generators. They reconstruct the stated first-order packet under the project hypotheses. Their converse does not construct a global Frey curve or all further integral lifts.
}{Group multiplication constrains a matrix packet more strongly than isolated coefficient values.}

\GuideSpread{Semilinearity retains a coefficient conjugation}{Frobenius on every element of $\F_{49}$}{10_semilinearity.pdf}{
The coefficient Frobenius is $\sigma(x)=x^7$. Since $w^2=5$ and $w^7=-w$, it sends $a+bw$ to $a-bw$. The guide checks this identity for all 49 elements.\par
The fixed line is $\F_7$. The negated line is $w\F_7$. These are one-dimensional spaces over $\F_7$, not two different coefficient fields.\par
The arithmetic semilinear statement also involves a specified lift of the base-field automorphism and one shared frame and reference. Under that additional pointwise condition, the remaining minus-congruence character coefficients lie in $w\F_7$. It reduces a conditional count of first-stage choices. It does not establish a nonempty affine solution set or a semilinear isomorphism of integral Tate lattices.
}{A coefficient identity and an arithmetic compatibility hypothesis should never be merged silently.}

\GuideSpread{The bounded support record is complete within its contract}{Four frozen supports with nine primes each}{11_support_sets.pdf}{
Every colored cell marks membership in one complete stipulated beta-support. The ten displayed rational primes form the union of the four nine-prime sets. The guide copies the source supports exactly and checks their cardinalities.\par
These patterns come from the depth-one, all-inert, plus-component necessary system with complete beta-support and every beta-ramifying rational prime at most 2000. The word complete refers to beta in that contract. It does not bound every divisor of $C$ or every upper correction.\par
The eight surviving coefficient tuples depend on inherited natural-model and trace assumptions. The support table is exact finite evidence for that necessary system. It does not classify every possible Beal counterexample, unrestricted support, lattice depth, or signature branch.
}{A finite support table is useful only with its completeness assumptions and exclusions attached.}

\GuideSpread{Eight residual tuples remain in the bounded branch}{Two coefficient-conjugate choices on each support}{12_eight_patterns.pdf}{
Patterns A and B retain $u=\pm3w$, while C and D retain $u=\pm w$. Every tuple has normalized ratio $r=6$ and coefficient $a_5=1$. Thus four supports yield eight normalized tuples.\par
These are survivors of the stipulated necessary equations and finite trace tests. They are not eight global representations or eight integer solutions. None is eliminated by the current affine-stage computation.\par
A common nonzero change-of-basis scalar is retained in the fixed-local comparison. Setting that scalar to one without justification can create a false exclusion. The research archive preserves the corrected calculation and its positive controls. Its earlier local tests leave all eight tuples, with their own allowed base residues.
}{A surviving necessary pattern is a candidate for further arithmetic analysis, rather than an existence theorem.}

\GuideSpread{The local Artin normalization changes the inertia value}{Exact conversion factors in $\F_7$}{13_artin_conversion.pdf}{
The frozen twisted residue generator and the ordinary-character inertia generator have different normalizations. Their conversion factor is displayed at each prime in the union support. A bar height denotes a residue in $\{0,\ldots,6\}$, not an ordered finite-field magnitude.\par
For $q\equiv1\pmod7$, conversion uses the norm from $\F_{q^2}$ to $\F_q$. For $q\equiv-1\pmod7$, it uses the direct seventh-power quotient. Multiplying the frozen inertia coefficient by this factor gives the ordinary Artin value $c_q$.\par
The affine toric row combines $c_q$ with the beta Frobenius value $b_q$. Omitting the scale would combine values in different local conventions and change the linear problem.
}{A normalization factor belongs to the arithmetic equation, even when it is a nonzero scalar.}

\GuideSpread{Ramified Frobenius needs its self-conductor term}{A norm-$q$ generator retains a nonzero local unit}{14_self_conductor.pdf}{
At a plus-congruence conductor prime $q$, the project supplies an integral cyclotomic generator $\pi_q$ with exact norm $q$. At its selected local place, write $\pi_q=q u_q$ with $u_q$ a unit. The uniformizer is the rational prime $q$.\par
The self-conductor contribution uses this unit and the other conjugate residue values. It cannot be replaced by an unramified good-prime test or a logarithm of zero.\par
Two arithmetic implementations check the norm and the self-term. The separate check lifts a cyclotomic root modulo $q^2$ and evaluates $u_q$ directly. Frobenius and inertia values must then be transferred to $K$ using the same Artin sign and degree conventions.
}{Evaluate a conductor place in its ramified local normalization.}

\GuideSpread{The bare first-stage equation has a global primitive}{Restricted existence under the frozen residual hypotheses}{15_bare_primitive.pdf}{
Let $\Sigma$ contain the complete beta-support, places above $2,3,5,7$, and the real places. The residual classes descend to this restricted Galois group. The stipulated local alpha class vanishes at beta-support, and the fixed-place cup classes vanish.\par
The written argument proves that the dual restricted localization kernel is zero, using class number one and the nonseventh local unit $\phi$. Finite-$\Sigma$ Poitou--Tate duality then gives $[\alpha\cup\beta]=0$.\par
A continuous normalized cochain $p_0$ therefore exists with $dp_0=-\alpha\cup\beta$, without auxiliary ramification. The result is an existence theorem. This edition supplies no defining field or numerical values for that one global primitive, and no assertion that it meets the required affine conditions.
}{The bare cup equation is solved, while its geometric affine origin remains unevaluated.}

\GuideSpread{The local target has twenty-four ordinary coordinates}{Three fixed, three seven-adic, and eighteen support coordinates}{16_local_coordinates.pdf}{
At places above $2,3,5$, the odd-primary ordinary-character spaces each have one unramified coordinate. At seven, the local basis uses $\phi$, $3$, and the uniformizer $7$.\par
Each of the nine beta-supporting places has two ordinary coordinates, labelled inertia and Frobenius. These are values on prescribed local Artin generators. Thus the local target dimension is $3+3+2\cdot9=24$ over $k$.\par
The picture numbers the nine support slots rather than repeating all prime labels. Their actual order is fixed separately for each pattern. These coordinates describe the ordinary characters $z=p-p_0$. The cochain $p$ itself need not be an ordinary character, since its differential is the prescribed cup product.
}{Apply character localization to differences of primitives, rather than to the primitive alone.}

\GuideSpread{Reciprocity computes the entire localization image}{A $14\times24$ relation matrix and a $24\times10$ image matrix}{17_global_localization.pdf}{
The fourteen global $\Sigma$-units are $\phi,2,3,\sqrt5,7$, and the nine support primes. Their principal-idele relations define a matrix $L$ with rank fourteen. Class number one supplies no further ideal-class variables.\par
The explicit matrix $J$ has ten independent columns, corresponding to the seven-adic ordinary unit direction and nine inertia directions. Exact checks give $LJ=0$, $\operatorname{rank}J=10$, and $\operatorname{im}J=\ker L$.\par
The arithmetic identification of this kernel with global localization uses the stated reciprocity theorem and hypotheses. It is stronger than a necessary finite overapproximation. It still classifies ordinary differences, rather than choosing the unknown origin of the affine primitive torsor.
}{A complete homogeneous image is not the same thing as an evaluated affine target.}

\GuideSpread{Projecting to the geometric conditions gives fifteen rows}{The exact Pattern C matrix $A=EJ$}{18_toric_projection.pdf}{
The local constraint map $E$ selects six fixed-local coordinates and nine toric combinations. Composing it with the global localization image gives $A=EJ$, a $15\times10$ matrix.\par
The figure shows the Pattern C matrix exactly. The copper divider separates the first six rows from the nine support rows. Color encodes residues zero through six in $\F_7$, with no ordering interpretation in the field.\par
Every toric row includes the correctly normalized Frobenius value at its own conductor prime. The guide independently multiplies $EJ$, verifies the stored matrix, and computes rank by exact elimination modulo seven. Extension of these matrices to $k$ preserves their ranks. Their target is the unknown geometric affine vector.
}{Compose the actual local condition with the complete global image before computing its obstruction.}

\GuideSpread{Kernel size and affine obstruction size are different}{Exact dimensions over $k=\F_{49}$}{19_affine_dimensions.pdf}{
The ranks for A, B, C, D are $7,7,8,7$. Since the domain has dimension ten, the homogeneous kernels have dimensions $3,3,2,3$. Since the target has dimension fifteen, the affine cokernels have dimensions $8,8,7,8$.\par
The kernel gives freedom among character corrections once an affine solution exists. The cokernel gives independent linear coordinates on the potential target obstruction. It does not imply that the actual arithmetic target can vary arbitrarily in those coordinates.\par
If the affine set is nonempty, the unrestricted first-stage counts are $49^3,49^3,49^2,49^3$. With the extra common pointwise semilinear condition, the corresponding bounds are $343,343,49,343$. Neither count constructs an arithmetic solution.
}{Compute a nonempty affine fiber before interpreting its homogeneous freedom as choices.}

\GuideSpread{A small minor explains Pattern C's improvement}{Two former plus directions are removed}{20_pattern_C_minor.pdf}{
Two old homogeneous plus directions of Pattern C are evaluated by the additional toric rows at 197 and 337. In the retained old basis, their matrix is $\bigl(\begin{smallmatrix}3&1\\5&2\end{smallmatrix}\bigr)$. Its determinant is one modulo seven.\par
Those two directions are therefore removed, reducing the old kernel dimension from four to two. The remaining ordinary-character directions are supported at 307 and 1483. For A, B, D, the remaining supports are 307, 797, and 1483.\par
This is a statement about differences of global primitives after all first-stage conditions. It does not force the absolute inertia of the unknown primitive $p_0$ to vanish at every other support place. It is a rank improvement, rather than a pattern elimination.
}{An invertible minor can remove homogeneous directions without deciding affine compatibility.}

\GuideSpread{The arithmetic question is now an affine cokernel test}{One global primitive determines one target vector}{21_affine_criterion.pdf}{
For a bare global $p_0$, the target has six fixed-local difference coordinates. At a support prime its entry is $-h_qc_q-c_qp_0(F_q)+b_qp_0(T_q)$. All values use the same frame, reference, and local generators.\par
A character correction $z$ satisfies every first-stage condition exactly when $Az=t(p_0)$. With a full left-null basis $N$, this is equivalent to $Nt(p_0)=0$. The obstruction has eight coordinates for A, B, D, and seven for C.\par
The matrices and target evaluator are ready. No actual arithmetic $t(p_0)$ or $Nt(p_0)$ has been supplied or computed. Substituting a synthetic zero vector would not evaluate this arithmetic problem.
}{The next step is an authenticated target calculation, rather than another homogeneous rank.}

\GuideSpread{Changing the bare primitive preserves the obstruction}{An exact synthetic origin-invariance control}{22_origin_invariance.pdf}{
Replacing $p_0$ by $p_0+z_{\mathrm{char}}$, where the global ordinary character $z_{\mathrm{char}}$ has parameter vector $z_0$, shifts the target to $t-Az_0$. Since $NA=0$, its obstruction coordinates remain $Nt$. The class in $\operatorname{coker}A$ is independent of the chosen bare primitive.\par
The displayed Pattern C control uses the synthetic target $t=e_1$ and correction $z_0=e_1$. Exact modular multiplication gives the same nonzero obstruction before and after translation.\par
This tests the evaluator and the linear origin-invariance identity. The target is explicitly synthetic. It is not obtained from one global arithmetic primitive, and its nonzero output does not eliminate an actual residual pattern. Other synthetic image controls similarly verify correction reconstruction without proving arithmetic existence.
}{Distinguish a canonical obstruction formula from an actual evaluation of its arithmetic class.}

\GuideSpread{The upper stage and integral endpoint are still separate}{Additional bad divisors may carry upper ramification}{23_upper_integral_gate.pdf}{
At a beta-unramified variable place, square-zero inertia over the field $k$ forces $p(T)=0$. It does not force $q(T)=0$. Thus the justified finite ramification set for Stage P does not bound the entire upper stage.\par
If upper first-order inertia is nonzero, the natural integral packet has content one. In the pinned frame, $N_{12}$ has valuation one, the diagonal entries are divisible by $49$, and $N_{21}$ is divisible by $343$. If the content is at least two, all first-order inertia vanishes modulo $49$.\par
These are local matrix conclusions with explicit higher-depth constraints. No global Frey converse or unrestricted upper conductor bound follows. Even a zero affine obstruction would leave these later realization problems.
}{Passing Stage P creates a new input for the upper and integral calculations.}

\GuideSpread{The next obligation is arithmetic, and its scope is precise}{Operation Komorebi's route remains incomplete}{24_research_boundary.pdf}{
Compute the canonical arithmetic class $Nt(p_0)$ for each remaining stipulated pattern. A nonzero result would exclude that first-stage branch. A zero result would give a finite family for the upper $q$ and integral compatibility problem.\par
Unrestricted support and depth, the minus component, other modular branches, and other signatures remain separate research obligations. The eight necessary patterns are not eliminated by the present release.\par
A September 2026 preprint by Peter Chocian (arXiv:2609.26996) claims a proof for primitive signature $(3,5,7)$. The monograph records that external claim and its source. Its full proof has not been audited in this edition. This guide states the status of the Operation Komorebi route, and makes no claim that Beal is solved.
}{Retain exact local progress and an explicit arithmetic boundary in the same research record.}
